\documentclass[11pt]{article}
\usepackage[T1]{fontenc}
\usepackage{lmodern,amsmath,amssymb,amsthm,booktabs,microtype}
\usepackage[margin=1in]{geometry}
\usepackage[hidelinks]{hyperref}
\usepackage{enumitem}
\setlist{nosep}
\newtheorem{theorem}{Theorem}
\newtheorem{corollary}[theorem]{Corollary}
\newtheorem{lemma}[theorem]{Lemma}
\allowdisplaybreaks
\setlength{\emergencystretch}{2em}
\title{A singular diagram and six bifurcations of planar Newtonian six-body central configurations}
\author{Chatgpt}
\date{Research manuscript, 27 November 2026}
\begin{document}
\maketitle
\begin{abstract}
We study planar Newtonian central configurations with mass vector
$(1,1,1,1,1,t)$ over a rational function field and at real positive parameters.
We rule out Tomar's diagram 98 on this line over the geometric generic
parameter, and at $t=3$. We extend the exclusion to a nonempty Zariski open subset of the three-parameter family $(1,1,1,a,b,c)$, including $(a,b,c)=(2,3,7)$. The original leading torque vanishes on an
equilateral factor. Expanding the normalized potential to cubic order
yields an additional necessary condition. A finite resultant witness
and exact interval calculations exclude both its ordinary and complex
collinear equal-triangle lifts. We also prove a lower bound of 4104
simple points in the unrestricted geometric generic fibre, assuming the
previous small-mass count. Endpoint catalogues and local certificates
identify three folds and three symmetry-breaking pitchforks in $0<t<1$.
They account for the observed net change of 480 labelled oriented
real configurations. Global real counts between the events and generic
finiteness of the entire six-body fibre on this line remain open.
\end{abstract}
\section{The mass parameter and the global target}
At a central configuration each gravitational acceleration is
$-\lambda(q_i-c)$, where $c$ is the centre of mass and $\lambda>0$.
We identify translations, rotations and positive dilations and retain all
labels unless otherwise specified. The preprint of Tomar~\cite{Tomar}
proves finiteness off an algebraic set containing this whole mass line.
Its approach studies a complex curve that would result from infinitely
many physical shapes at fixed masses. Its possible singular places have
combinatorial types called diagrams. Diagram 98 involves two triples
of bodies and a necessary leading torque balance between them.
Here $t$ denotes a mass parameter. It is distinct from the two projected
shape variables of the triples. In the section on diagram 98 the same
mass parameter is denoted by $k$.

\subsection*{What the full algebraic system already implies}
Translate the marked sixth body to $(0,0)$ and place body five at $(1,0)$
by the unique orientation-preserving similarity with these properties.
The four remaining positions give eight real coordinates $x$.
Adjoin inverse distances $s_{ij}$ by
$s_{ij}^2\lVert q_i-q_j\rVert^2=1$, using the positive sheet for real
physical solutions and fixing $s_{56}=1$. With masses $m_i$, define
$U=\sum_{i<j}m_im_js_{ij}$ and
$I=(\sum m_i)^{-1}\sum_{i<j}m_im_j\lVert q_i-q_j\rVert^2$.
The shape equations are the eight derivatives of $U\sqrt I$ in this
chart. They have a rational formulation
$2I\,\partial_x U+U\,\partial_x I=0$, using implicit differentiation
on the distance sheets. Restrict to $I(5+t)\ne0$. All physical solutions
in the positive mass range are represented. The distance equations are
\'{e}tale in their auxiliary variables at physical solutions, so a
nondegenerate physical shape gives a simple zero of this square algebraic
system.

\begin{theorem}
Over $\overline{\mathbb Q(t)}$ this normalized system has at least 4104
distinct simple solutions. If the whole coordinate algebra is finite over
$\mathbb Q(t)$, its dimension is at least 4104.
\end{theorem}
\begin{proof}
The small-mass theorem in~\cite{Small} gives 4104 distinct nondegenerate
labelled oriented solutions for every sufficiently small positive $t$.
Over an algebraically closed field the existence of at least 4104
distinct zeros with nonzero Jacobian determinant is a constructible
condition on $t$. Indeed it is the image of the locus of 4104 ordered
tuples of such solutions, deleting the pairwise diagonals. Its image in
the affine line contains infinitely many real points, hence a nonempty
Zariski open subset. This proves the assertion at the geometric generic
point. Each simple point contributes one to the dimension when the
algebra is finite. This argument permits other positive-dimensional
components.
\end{proof}
For a proof of the desired upper finiteness bound, each shape coordinate
must satisfy a monic polynomial over $\mathbb Q(t)$ on every component
under consideration, or all remaining components must be shown unable
to contain a physical curve. A sufficient obstruction for one diagram
alone does not supply this conclusion.

\section{A reduction of the diagram-98 torque condition on the five-equal-mass line}

Consider masses proportional to \( (1,1,1,1,1,k) \), with \(k>0\). In a diagram-98 singular place, two triples separate in opposite complex coordinates. After interchanging the triples and relabelling bodies, their masses are \(\mu=(1,1,1)\) and \(\nu=(1,1,k)\). This note restricts their leading coefficients. The cubic-order argument in the next section will exclude the remaining leading shapes at generic mass.

\textbf{Proposition.} There exists an explicitly defined nonzero polynomial \(P\in\mathbb Z[k]\) with \(P(3)\ne0\) such that, for every \(k>0\) with \(P(k)\ne0\), every solution to the diagram-98 leading shape and cross-torque conditions has both of the following properties.

\begin{enumerate}
\item The projected shape \(x\) of the all-equal triple satisfies \(x^2-x+1=0\).
\item The cross sign matrix is of rank one over the complex numbers. Its signs can be written as \(\epsilon_{jl}=\delta_j\eta_l\), independently of the chosen square-root branches.
\end{enumerate}

The proposition includes \(k=3\). It covers every labelled allocation of the five interchangeable unit masses. The projected complex shape need not represent a real equilateral triangle.

\subsection*{Input from the central-configuration equations}

Tomar, \emph{Smale's sixth problem for six bodies in the plane: finiteness of central configurations for generic masses, with a proof formalized in Lean} (Hexagon 2611.00377v1, Sections~6.4--6.6), Lemmas 6.9 and 6.11, and Theorem 6.12, gives the following necessary conditions. For the triple with masses \(m=(1,1,h)\) and shape \(z\), set

\[
(\ell_1,\ell_2,\ell_3)=(1+hz,hz-1-h,1-2z),\qquad
\Lambda_j=\prod_{i\ne j}\ell_i.
\]

All three linear forms are nonzero at a diagram-98 place. Write \(S_h\) for Tomar's degree-34 shape polynomial. With \(S_1(x)=S_k(y)=0\), the cross-torque vanishes only if some signed sum

\[
\sum_{j,l=1}^{3}\epsilon_{jl} r_j s_l = 0,\qquad
r_j^2=\mu_j^2\Lambda_j(x;1),\quad
s_l^2=\nu_l^2\Lambda_l(y;k)
\]

vanishes. The signs of a row or column change when a corresponding square-root branch changes. Rank and the equation with the transformed sign matrix are invariant. Overall sign is irrelevant, so require \(\epsilon_{11}=1\). Tomar's Remark 6.14 observes that rank-one matrices include the factorable sums that vanish for an equal-mass triple. The corresponding original obstruction polynomial therefore vanishes identically on this line. The proposition tests the other signs and the other equal-triple shapes separately.

\subsection*{Reduction of the shape equations}

Put \(E(z)=z^2-z+1\) and

\[
L=E(z)(z^4-5z^3+9z^2-5z+1)(z^4+z^3-2z+1)(z^4-2z^3+z+1).
\]

For the equal triple retain the degree-26 polynomial

\[
\begin{aligned}A=L{}&(2z^4-z^3+2z^2-2z+1)(z^4-2z^3+2z^2-z+2)\\
&(2z^4-7z^3+11z^2-7z+2).\end{aligned}
\]

For the other triple put

\[
\begin{aligned}
B_k=L{}&[2z^5-5z^4+4z^3-(3k+3)z^2+(3k+2)z-(k+1)]\\
&[z^4-2z^3+(k+1)z^2-kz+k+1]\\
&[2z^5-5z^4+4z^3+(3k+1)z^2-(3k+2)z+k+1]\,
[z^2-z-k-1].
\end{aligned}
\]

Its degree is 30 for every \(k\). Direct multiplication in \(\mathbb Z[k,z]\) gives

\[
\begin{aligned}S_1&=A(z)(2z-1)^2 E(z)(z+1)^2(z-2)^2,\\
S_k&=B_k(z)(2z-1)^2 E(z).\end{aligned}
\]

In both products the retained polynomial already contains \(E\). Every other removed root makes at least one \(\ell_i\) vanish. Thus these equations retain every admissible shape even at special masses. The equal-triple polynomial \(A/E\) has no common root with \(E\); in particular, a non-equilateral equal-triple branch cannot specialize to an equilateral root at 3. Both retained polynomials have constant leading coefficient 4; no shape escapes to infinity by a degree drop. These identities and the checks on specialization are reproduced by \texttt{SourceAlgebra.py}.

\subsection*{An integer obstruction polynomial}

Let \(\mathcal N\) be the 240 canonical matrices of rank greater than one and \(\mathcal F\) the 16 canonical matrices of rank one. For either class \(\mathcal C\), define

\[
H_{\mathcal C}(r,s)=
\prod_{e\in\mathcal C}\left(\sum_{j,l} e_{jl}r_js_l\right).
\]

Changing the sign of any \(r_j\) or \(s_l\) permutes these factors up to overall signs. If \(e_{11}\) is affected, \emph{every} factor needs a minus sign to restore the normalization, and the number of factors is even. Hence every exponent of every radical in the product is even. Substituting the six square identities above yields \(h_{\mathcal C}(x,y,k)\in\mathbb Z[x,y,k]\). Define

\[
P(k)=\operatorname{Res}_x\big(A,\operatorname{Res}_y(B_k,h_{\mathcal N})\big)
\operatorname{Res}_x\big(A/E,\operatorname{Res}_y(B_k,h_{\mathcal F})\big).
\]

All polynomial evaluations take place after replacing radicals by their squares in the invariant product. This is an explicit polynomial specified by products and determinants rather than by expanded coefficients. A root pair with either a nonfactorable sign matrix, or a non-equilateral equal-triple shape and a factorable matrix, makes the corresponding nested resultant zero.

It suffices to check \(P(3)\ne0\). At 3 both shape polynomials have their full degrees and all roots are simple. The product formula for a resultant of a fixed-degree outer polynomial multiplies the inner values over all outer roots, up to a nonzero power of the leading coefficient, even if the degree of the inner polynomial drops in a specialization. Thus the finite check below establishes \(P(3)\ne0\) without assuming that a component of any system of rational equations meets the tested affine chart. No inversion of a parameter, radical, or shape coordinate is needed to construct \(P\). For a sign class of size \(M\), \(H_{\mathcal C}\) is homogeneous of degree \(M/2\) in the squared coordinates of each triple separately. The substituted entries have shape degree at most two, and the mixed entries have parameter degree at most four. Hence \(h_{\mathcal C}\) has degrees at most \((M,M,2M)\) in \((x,y,k)\). The mixed-shape polynomial has parameter degree at most four and shape degree 30. The Sylvester determinant bounds the inner resultant's parameter degree by \(30(2M)+4M=64M\); the outer shape polynomials are constant in the parameter. In particular,

\[
\deg P \le 26\cdot64\cdot240+24\cdot64\cdot16=423936.
\]

The finite calculation below proves nonvanishing, and hence the finite-exception assertion.

\subsection*{The finite witness}

The computation at 3 uses \(a_j=\mu_j/\sqrt{\ell_j(x;1)}\) and \(b_l=\nu_l/\sqrt{\ell_l(y;3)}\). All linear forms are nonzero on both root sets, so multiplying by one nonzero common scalar for each triangle, with harmless changes of row/column signs, identifies this test with the signed sums in \(r,s\).

The input file \texttt{seeds98.json} contains Gaussian dyadic centres of denominator \(B=2^{48}\). A floating-point calculation can be used to generate them; the verifier treats them as arbitrary untrusted input. Set \(R=2^{-39}\) and \(H=2^{-32}\). For a proposed root centre \(c\), write the exact Taylor coefficients of its assigned factor as \(f(c+w)=\sum d_i w^i\). The verifier proves \(|d_1|R>|d_0|+\sum_{i\ge2}|d_i|R^i\). Rouché's theorem gives exactly one zero in each disk. Disjointness and the degrees certify completeness.

For the proposed radical centre \(w\), let \(\ell_0=\ell(c)\) and let \(s\) be the slope of this linear form. The true value differs from \(\ell_0\) by at most \(|s|R\). The verifier proves

\[
2|\ell_0w|H>
|\ell_0w^2-m_j^2|+|\ell_0|H^2+|s|R(|w|+H)^2.
\]

A second application of Rouché's theorem isolates a radical root in each disk. The other root is its negative. The global sign enumeration covers all choices. Given centres \(w_j,v_l\) with error at most \(H\), the product \(a_jb_l\) has error at most \((|w_j|+|v_l|)H+H^2\). The verifier sums these errors and compares them to a lower bound on the modulus of every signed sum. It bounds \(|z|\) below by \(\max(|\Re z|,|\Im z|)\) and above by \(|\Re z|+|\Im z|\), so every comparison uses integers or exact rational arithmetic.

There are 26 equal-triple roots and 30 mixed-triple roots. All \(26\cdot30\cdot240=187200\) rank-greater-than-one combinations are excluded. All \(24\cdot30\cdot16=11520\) non-equilateral rank-one combinations are also excluded. The least certified lower bound on the modulus over these two classes exceeds 0.036477197 in the inverse-root coordinates used for this finite witness. Among 199680 total combinations, 240 fall within the interval error. Every one has an equilateral equal-triple root and factorable signs. This count is a count of signed algebraic tests, not of central configurations.

For either equilateral root, set \(Y_j=\mu_j^2\Lambda_j(x;1)\). The product of the four three-term signed sums equals

\[
Y_1^2+Y_2^2+Y_3^2-2(Y_1Y_2+Y_2Y_3+Y_3Y_1)=9E(x)^2.
\]

It follows that for both equilateral roots some signed sum over the equal triangle vanishes. Multiplying this identity by an arbitrary signed sum over the mixed triangle produces a factorable solution of the scalar torque equation. Additional information from the governing equations is required to decide whether any such leading solution extends to a singular place.

\subsection*{A force identity on the surviving branch}

There is a useful further restriction in the physical leading coordinates. Suppose
\(k>0\), \(k\ne1\), and \(P(k)\ne0\). Write \(Z_j\) for the leading
projected coordinate of equal body \(j\), and \(W'_l\) for that of mixed
body \(l\), in the two opposite coordinates of Tomar's Lemma 6.11. Then
there are nonzero \(\gamma_l\) such that for the leading cross inverse
distances
\[
   \widehat p_{jl}=\gamma_l Z_j .
\]
In particular, for each mixed body \(l\), the three values
\(\widehat p_{jl}^3\) are equal, and
\[
\sum_{j=1}^3\widehat p_{jl}=0,\qquad
\sum_{j=1}^3\widehat p_{jl}^3 Z_j=0.
\]

To see this, first rule out a cancellation in the mixed triple.
At its shape \(y\), set \(Y_l=\nu_l^2\Lambda_l(y;k)\) and
\(D=\sum_lY_l^2-2\sum_{l<q}Y_lY_q\). Calculating the resultant with
each of the eight displayed factors of \(B_k\) and applying Sturm's
algorithm to its irreducible factors finds just one positive root in
their union, namely \(k=1\). The exact calculation is in
\texttt{SourceAlgebra.py}. Thus the mixed factor of a rank-one
vanishing sum is nonzero, so the equal factor must vanish.

The three equal masses have centered coordinates proportional, in some
order, to \((1,\omega,\omega^2)\), where \(\omega\) is a primitive cube
root of unity. This follows by substituting \(x^2-x+1=0\) into the
linear forms above. The three numbers \(Z_j^3\) have a common nonzero
value, so for a common constant \(c\) one possible set of radicals
of \(1/Z_j\) is \(cZ_j\). These sum to zero. A signed zero sum of
three nonzero complex numbers has only one sign pattern up to a
global sign: two inequivalent patterns would isolate a single nonzero
term after adding or subtracting the relations. The equal-triple
radicals in a vanishing factor must therefore be proportional to
\(Z_j\). Factorability and
\(\widehat p_{jl}^{\,2}Z_jW'_l=-1\) give the displayed formula.
Both moment identities follow from \(\sum_j Z_j=0\).

\subsection*{Reproduction}

Run \texttt{python3 Diagram98.py} (Python 3.10+ standard library) and \texttt{python3 SourceAlgebra.py} (SymPy). The optional \texttt{generate98.py} regenerates the untrusted dyadic input with NumPy. The program certifies the finite specialization with exact arithmetic. The specialization and elimination argument is the mathematical proof above.

\section{Cubic-order exclusion of diagram 98}

\textbf{Proposition.} Outside a finite algebraic set of parameters
\(k>0\), no diagram-98 singular place exists for the masses
\((1,1,1,1,1,k)\). The exclusion holds at \(k=3\).

\paragraph{A chart near the two triangles.}
Use Tomar's exponents in Lemma 6.10. By a complex rotation fix
\(z^A_{12}=1\), and set \(h^2=1/w^B_{12}\) in these rotated
coordinates. A further quadratic extension to extract \(h\) is harmless.
Write the physical positions in the chart as
\[
 (Z_j,W_j)\quad(j\in A),\qquad (h^2 Z'_l,h^{-2} W'_l)\quad(l\in B).
\]
Thus \(Z_{12}=W'_{12}=1\), \(h\to0\), and all the coordinate
vectors in capital letters have finite limits. The designated pair
differences and all the individual \(Z_j,W'_l\) at the limit are nonzero.
The two triples obey their own isolated normalized three-body equations at
the limit. The two triple centres are included as chart variables; stationarity of the harmonic translation terms recovers the centre-of-mass condition. For \(h\ne0\),
the chart gives configurations modulo the common rotation, locally up
to the choice of sign of \(h\).

With each local inverse-distance branch chosen by the given place, the
normalized central equations imply stationarity of
\(V=\frac12\sum_i m_i z_iw_i+\sum_{i<j}m_im_jp_{ij}\).
Write \(V_A,V_B\) for the analogous isolated-triple functions. Each cross
inverse distance, divided by \(h\), is locally analytic and equals
\[
 (-Z_jW'_l)^{-1/2}
 (1-h^2Z'_l/Z_j)^{-1/2}(1-h^2W_j/W'_l)^{-1/2},
\]
with its own leading sign. Thus, as an analytic function of all chart
coordinates,
\[
 V=V_A+V_B+h C_0+h^3 C_1+O(h^5).
\]
The preceding section gives a rank-one sign matrix, hence
\(C_0=a(Z)b(W')\), where the signed terms of \(a\) are branches of
\(1/\sqrt{Z_j}\) and those of \(b\) are branches of a common nonzero
constant times \(\nu_l/\sqrt{W'_l}\). At the limiting equal triple
\(a=0\), the signed radicals in \(a\) are \(c Z_j\), and
\(\partial a/\partial Z_j=-c/2\) for all three bodies. The mixed
triple has \(b\ne0\) for \(k>0\), \(k\ne1\), by the signed-cancellation
test in the preceding section. Denote by \(g_{jl}\) the leading cross
radical at the limiting chart point. Then
\[
 C_1=\frac12\sum_{j,l}\nu_l g_{jl}
    \left(\frac{Z'_l}{Z_j}+\frac{W_j}{W'_l}\right)
 \quad\text{at that point}.
\]

\paragraph{Possible equal-triple lifts.}
Put \(x=Z_{13}\), \(v=W_{12}\), and \(x'=W_{13}/v\).
At the limit \(x=r\), where \(r^2-r+1=0\).
For a Lagrange-type solution of the three-body equations,
all three internal inverse distances have cube \(1/3\).
The ratios \(d=D_{13}/D_{12}\) and \(e=D_{23}/D_{12}\), with
\(D_{ij}=Z_{ij}W_{ij}\), satisfy \(d^3=e^3=1\),
\(r^2-(1+d-e)r+d=0\), and \(x'=d/r\).
Checking the three cube roots gives either \(d=e=1\),
\(x'=1/r\), or \(d=r^2,e=(r-1)^2\), \(x'=r\).
For an Euler-type lift, \(x'=r\). Here the distances are
\(v,vr^2,v(r-1)^2\), which all have the same cube. The three inverse
cubes therefore equal \(c\epsilon_{ij}\) for a common scalar \(c\ne0\)
and real signs \(\epsilon_{ij}\). The weighted graph Laplacian on the
three equal bodies with these edge signs is real symmetric, and the
centered vector \(Z\) is its eigenvector with eigenvalue \(1/c\).
This eigenvalue is real. Its conjugate vector is a second, independent
centered eigenvector, since the ratio \(Z_{13}/Z_{12}=r\) is nonreal.
The Laplacian is therefore scalar on the centered plane and zero on
the constant vector. Its three off-diagonal entries are equal, so all
three inverse cubes agree and the force equation gives their common
value \(1/3\).
Thus \(p_{ij}^3=1/3\) at both kinds of lift. Let
\(I_A=\sum_j Z_j W_j=(D_{12}+D_{13}+D_{23})/3\) at the centered
limit. If \(x'=1/r\), all three \(D_{ij}\) are equal and \(I_A^3=9\).
If \(x'=r\), \(I_A=0\) because the three \(D_{ij}\) sum to zero.

\paragraph{Regular equal-triple variables.}
In the independent centre variables \(z_c,w_c\), the translation part
of \(V_A\) is \(3z_cw_c/2\), with invertible Hessian. The centered
part is \(\sum_{ij} f_{ij}(D_{ij})\), where
\(f_{ij}(D)=D/6+D^{-1/2}\) with the chosen inverse-distance signs.
At both lifts \(f_{ij}'=0\) and \(f_{ij}''=1/(4D_{ij})\).
Indeed \(p^2D=1\) and \(p^3=1/3\).
The three distance variables are
\[
  (D_{12},D_{13},D_{23})=(v,vxx',v(x-1)(x'-1)).
\]
If \(x'\ne x\), their Jacobian in \((v,x,x')\) is invertible; hence
all five equal-triple coordinates in the rotation slice are regular.

At the collinear lift set \(p=(x+x')/2\), \(u=(x-x')/2\).
Here the last two distances are \(v(p^2-u^2)\) and
\(v((p-1)^2-u^2)\). In the \((v,p)\) Hessian block,
\(H_{vv}=\sum D_{ij}/(4v^2)=0\) and
\(H_{vp}=\sum(\partial_pD_{ij})/(4v)=(2r-1)/2\ne0\).
Thus the four variables \(z_c,w_c,v,p\) are regular even at this lift.
At \(h=0\) eliminating them gives analytic functions of \(u^2\), with
centres zero and \(p=r+O(u^2)\). All distance variations in this
solution are \(O(u^2)\). Since all \(f'_{ij}=0\), the reduced
equal-triple potential is constant up to \(O(u^4)\).

\paragraph{The critical value in the remaining variables.}
Treat the mixed-triple coordinates, denoted by \(Y\), as independent
parameters in a neighbourhood of their limiting values \(Y_0\).
Temporarily truncate the potential at
\(V^{\mathrm{tr}}=V_A+V_B+h a(Z)b(W')\). For every \(h,Y\) near
the limiting point, the limiting equal triangle with just its three
\(W_j\) translated by the same quantity
\(-2h b(W')\partial_{Z_1}a\) is exactly stationary in \emph{all}
equal-triple variables of the truncated problem. In fact the
translation leaves internal distances and all \(Z_j\) unchanged,
and the change of each \(Z_j\)-derivative of the harmonic term is half
the common translation. The value of \(V_A\) is unchanged since
\(\sum_j Z_j=0\). This observation holds in moving as well as fixed
mixed-triangle centres.

Use the implicit function theorem on the regular equal-triple variables
just identified, first for the truncation, then for \(V\). The full and
truncated critical variables differ by an analytic multiple of \(h^3\).
Thus the two reduced potentials differ by \(h^3 J\), where \(J\)
is analytic in \(h,Y\) and, in the collinear case, also \(u\).
At the limiting configuration \(J=C_1\). For example, this identity
follows by expanding the truncation about its stationary regular
variables, which gives only a quadratic error in their shift. This
argument can equally be applied in completed local power-series rings;
along a Puiseux branch the estimates below are estimates in its
uniformizing parameter. No Hessian of \(V_B\) is inverted.

In the noncollinear case the truncated critical value is
\(V_B(Y)+\mathrm{const}\) by the exact translation branch. The
partial derivative of the full critical value in \(h\), at fixed
\(Y\), consequently forces \(C_1=0\) in the limit.

In the collinear case the exact translation branch lies at \(u=0\);
it is critical also in \(u\) there. The reduced truncation therefore
has the analytic form
\[
  V_B(Y)+\mathrm{const}+u^2 L(h,u,Y),\qquad L(0,u,Y)=O(u^2).
\]
The coefficient \(d(Y)=\partial_h L(0,0,Y)\) is nonzero near
\(Y_0\). Indeed, differentiating the reduced truncation at \(h=0\)
gives precisely \(b(W')a(Z)\) at the regular critical variables there.
The centered \(Z\)-coordinates have shape \(x=p+u=r+u+O(u^2)\).
In these coordinates the vanishing equal-triple radical sum has an
\emph{exact double zero} at \(x=r\). The product of its four signed
choices is a nonvanishing analytic factor times \(9(x^2-x+1)^2\),
and the other three signed choices are nonzero by uniqueness of the
sign relation. This proves the claimed nonvanishing, since \(b(W')\ne0\).

Stationarity of the full reduced potential in \(h\) gives
\(0=u^2[d(Y)+O(h,u)]+O(h^2)\), hence \(u=O(h)\). Stationarity
in \(u\) then gives
\[
 0=2h d(Y)u+O(u^3+h u^2+h^2u+h^3),
\]
so \(u=O(h^2)\). Dividing the original \(h\)-derivative by \(h^2\)
now yields \(3J(0,0,Y_0)=0\), and again \(C_1=0\). These differentiations
are partial derivatives in the independent coordinates, using
stationarity, and do not require a rate estimate on \(Y-Y_0\).

\paragraph{The finite cubic torque test.}
The force identity in the preceding section gives \(g_{jl}=c Z_j b_l/\nu_l\),
where \(b_l\) is a signed choice of \(\nu_l/\sqrt{W'_l}\) and \(c\)
is nonzero. Thus the necessary equation is
\[
  3\sum_l b_l Z'_l + I_A\sum_l b_l/W'_l=0.
\]
Write \(M=2+k\) and
\(W(t)=(1+kt,kt-1-k,1-2t)/M\).
Then \(W'=W(y)\), where \(B_k(y)=0\). By Tomar's three-body
dichotomy, the opposite coordinates have one of the following forms.
In the Lagrange case \(Z'=R W(d/y)\), with \(d^3=1\) and \(R^3=M^2\).
This follows from the three distance equations and \(p_{ij}^3=1/M\).
In the Euler case \(Z'=\lambda W(y)\) where
\[
 \lambda^3=F^2,\qquad
 F=2+k\left(\sigma_{13}/y^2-\sigma_{23}/(y-1)^2\right).
\]
Indeed \(p_{ij}=\kappa\sigma_{ij}/W'_{ij}\), \(\kappa^2=1/\lambda\)
and the difference of the first two normalized force equations yields
\(\kappa^3F=1\).

In the Lagrange case set \(U=\sum b_l W_l(d/y)\), \(s=M^2\); in the
Euler case set \(U=\sum b_l W_l(y)\), \(s=F^2\). In both cases put
\(T=\sum b_l/W'_l\).
For the noncollinear equal lift the necessary test is \(3s U^3+T^3=0\);
for the collinear equal lift it is \(U=0\).
In either test the verification may replace \(b_l\) by a signed
\(\nu_l/\sqrt{\ell_l(y;k)}\), removing one common nonzero scalar.

At \(k=3\), \texttt{NextTorque.py} evaluates interval enclosures for
all 30 roots, all four overall-normalized radical sign choices, each
of three possible \(d\), and each of four possible Euler sign pairs.
It uses the previously certified disks inside dyadic rectangles, with
all operations rounded outwards on the integer grid \(2^{-80}\).
It includes combinations of shape and triangle type even when those
combinations do not satisfy their corresponding internal equations.
The 840 evaluations in each test give uniform lower bounds
\(2.721333\) for \(\lvert 3sU^3+T^3\rvert\), and \(0.050183\) for
\(\lvert U\rvert\). Thus neither necessary test can hold at \(3\).

For generic parameters, take in each test the product over the four
normalized sign choices. It is invariant under every radical sign
flip, since a global sign changes each odd factor's sign four times.
It is therefore a polynomial in the three radical squares, with rational
function coefficients in \(y,k,d\). Substitute
\(b_l^2=\nu_l^2/\ell_l(y;k)\), clear powers of
\(M y(y-1)\prod_l\ell_l\), and take the product over the three
Lagrange values \(d\), or respectively the four Euler sign pairs.
The resulting numerator polynomials have rational coefficients.
A common denominator may be used throughout. At \(k=3\) no one
of these denominators vanishes on the 30 shape roots: the previous
certificate excludes all \(\ell_l=0\), and \(B_3(0)B_3(1)\ne0\)
by direct evaluation. Their resultants in \(y\) against \(B_k\)
are therefore nonzero at \(k=3\) and vanish at any admissible
specialization exhibiting one of the tests. Multiplying these four
resultants, the polynomial of the previous section, and the factor
\(k-1\) gives a nonzero obstruction at 3. This proves the proposition.
The remaining diagram classes in the six-body problem require separate
arguments.

\section{Three equal masses and three independently variable masses}

Normalize the masses to \((1,1,1,a,b,c)\) with \(a,b,c>0\).
For a partition in diagram 98 that groups the three specified equal bodies
together, the exclusion just proved extends to all points outside a
proper algebraic subset of this three-parameter space.

Indeed, keep the same degree-26 polynomial for the equal triangle. For
\(m=(a,b,c)\), put \(M=a+b+c\) and
\[
 W(y)=(b+cy,cy-a-c,b-(a+b)y)/M.
\]
In Tomar's shape polynomial the Euler factor with signs \((1,-1,1)\)
is divisible identically by \(y^2-y+1\), which also occurs among the
Lagrange factors. Removing its duplicate gives a polynomial \(G_{a,b,c}\)
of degree 32 with leading coefficient \((a+b)^4\). This identity and
the specialization below are checked by \texttt{SourceAligned.py}.
The signed cross-torque resultants and the two cubic-order resultants
are constructed in the same way as for the line. In the Euler formula
replace the expression for \(F\) by
\[
 F=a+b+c\left(\sigma_{13}/y^2-\sigma_{23}/(y-1)^2\right).
\]
The mixed radical sum also enters the collinear argument. Its nonvanishing
at all mixed shapes on a Zariski-open set follows by taking a further
resultant of \(G_{a,b,c}\) with the four-sign product of that sum,
after clearing the same denominators. The reduction of the equal-triple
critical equations treats the mixed coordinates and their masses as
parameters throughout.

At \((a,b,c)=(1,2,7)\), the 32 mixed shapes are simple, and all
coordinates and denominators needed in the tests are nonzero.
\texttt{AlignedWitness.py} excludes the
\(26\cdot32\cdot240+24\cdot32\cdot16=211968\) leading combinations
specified in the proof, with a certified lower bound of \(0.007125\)
in the inverse-root gauge. For each equal lift it checks
\(32\cdot4\cdot7=896\) higher-order tests; the lower bounds for the
noncollinear test, collinear test and three-term mixed radical sum are
respectively \(23.192197\), \(0.113938\) and \(2.141114\).
In particular, each of the finitely many parameter polynomials used
in this aligned partition is nonzero in \(\mathbb Q[a,b,c]\).
The additional equality \(a=1\) at this witness is compatible with
the full-degree and denominator checks.

The same verifier checks a second aligned specialization, \((a,b,c)=(2,3,7)\),
where the distinguished triple comprises exactly the three unit masses.
Here the 211968 leading tests give lower bound \(0.031555\),
the 896 evaluations per cubic-order test give lower bounds \(7.729953\)
and \(0.065741\), and the mixed radical sums have lower bound
\(1.654861\). All 32 mixed shapes are simple and admissible for the
required divisions. Together with the three split-partition
checks at these same masses, this supplies one concrete point of the
common open set, \((1,1,1,2,3,7)\).

\subsection{Three designated equal masses distributed between the triangles}
Consider masses $(1,1,1,a,b,c)$ with $a,b,c>0$.
If one triangle contains precisely two of the three designated unit bodies,
rename the remaining masses so that the triangles are
$A=(1,1,a)$ and $B=(1,b,c)$.
Let $E(x)=x^2-x+1$, and let $S_m$ and $\ell_j(x;m)$ be the
polynomials in Tomar, Section~6.4. Set
\[
 A_a(x)=\frac{S_{(1,1,a)}(x)}{(2x-1)^2E(x)},\qquad
 B_{b,c}(y)=\frac{S_{(1,b,c)}(y)}{E(y)}.
\]
Both quotients are integral polynomials, of generic degrees $30,32$
in their respective shape variables with leading coefficients
$4,(1+b)^4$. Multiplication of the Lagrange and Euler factors verifies
these exact divisions. Both quotients retain a factor $E$.
The only other discarded root, $x=1/2$, satisfies
$\ell_3(x;1,1,a)=0$, which Tomar's Lemma~6.11 forbids.
Consequently all admissible projected shapes are roots of the displayed
polynomials at every positive parameter value.

Use the signed product over all 256 sign matrices with first entry $+1$,
as in Tomar, Theorem~6.12. Its factors are
$\sum_{j,l}e_{jl}r_j s_l$, where
$r_j^2=m_{A,j}^2\prod_{i\ne j}\ell_i(x;A)$ and
$s_l^2=m_{B,l}^2\prod_{i\ne l}\ell_i(y;B)$. The product
is invariant under each one of the six radical sign changes, hence yields
a polynomial $H(x,y;a,b,c)$ in integer coefficients by substitution.
An admissible leading solution makes $H$ vanish. Define
\[
 G(a,b,c)=\operatorname{Res}_x\!\left(A_a(x),
             \operatorname{Res}_y\big(B_{b,c}(y),H(x,y;a,b,c)\big)\right).
\]
Any diagram-98 place with the stated allocation has $G(a,b,c)=0$.

At $(a,b,c)=(3,2,7)$, both retained polynomials are square-free and
coprime to the products of their respective linear forms. A rational
interval certificate checks each of the $30\cdot32\cdot256=245760$
possible signed sums in the coordinates
$m_{A,j}/\sqrt{\ell_j(x;A)}$, $m_{B,l}/\sqrt{\ell_l(y;B)}$.
Use dyadic Gaussian centres with denominator $2^{48}$, disks of radii
$2^{-39}$ for shape roots and $2^{-32}$ for radicals, and the two
Rouch\'e comparisons in the preceding diagram-98 calculation.
The least certified lower bound on the modulus exceeds $2.28031515$.
Changing to the radicals in $H$ multiplies each triangle by a common
nonzero scalar and may change row or column signs, all of which are
included. Thus $G(3,2,7)\ne0$ by the product formula for the two
resultants. Positive specialization preserves the leading coefficients of
$A_a$ and $B_{b,c}$, including at the witness. The certificate
\texttt{Split98.py} also verifies the two symbolic divisions and their
specialization, disk coverage, all signed sums and their positive margins.

For prescribed labels fix the order of the two non-unit-labelled bodies in
the second triangle. Up to permutations of the three designated unit
bodies and exchange of the two triangles, the three split obstructions
are $G(a,b,c),G(b,a,c),G(c,a,b)$. At the common mass vector
$(1,1,1,2,3,7)$ the same procedure checks all three, using
$(2,3,7),(3,2,7),(7,2,3)$, respectively. Each check excludes all
245760 signed sums. The respective certified lower bounds on their
moduli exceed $0.23389010$, $2.28031515$, and $1.67447040$.
The products are nonzero there, as is the aligned-triangle obstruction
by the preceding calculation at mixed masses $(2,3,7)$. The product
of these four obstructions is therefore nonzero at $(a,b,c)=(2,3,7)$.
Every diagram-98 place is excluded on a Zariski-open subset of the
three-equal-mass locus containing this point and, in particular, on a
relative ordinary open neighbourhood where the three remaining masses
may vary independently.

\section{Endpoint symmetries and a forced two-axial degeneracy}
Single out body 6 even when its mass is one. In this section a marked
shape identifies permutations of bodies 1--5 and all orthogonal
transformations. A mirror-symmetric shape is classified by the number of
bodies on its axis. A nonidentity orientation-preserving symmetry must fix
the marked body. Its remaining five vertices form orbits of its order about
that point, since none can coincide with it. The order must therefore be
five. Its sole possible nontrivial rotational class is a regular pentagon
with the marked body at the centre. Every other symmetric shape has just
one reflection as a marked shape. A shape with no such reflection is
called chiral; the table counts a chiral shape and its mirror together.

\begin{center}
\begin{tabular}{lrrr}
\toprule
Marked shape type & Small positive $t$ & $t=1$ & Labelled oriented weight\\
\midrule
Collinear                    &3&3&120\\
Four bodies on its unique mirror &2&2&120\\
Two bodies on its unique mirror  &13&11&120\\
Regular pentagon plus centre     &1&1&24\\
Chiral pair                      &8&7&240\\
\midrule
Total marked shapes              &27&24&\\
Total labelled oriented shapes   &4104&3624&\\
\bottomrule
\end{tabular}
\end{center}

Here the inputs are the complete five- and six-body catalogues of
Moczurad and Zgliczy\'nski~\cite{MZ}, the restricted equilibria,
interval enclosures and nondegenerate continuation in~\cite{Small}, and
the released six-body enclosures. The accompanying
\texttt{symmetry\_counts.py} refines these data. At the light end it
enumerates the orthogonal group of each of the five primary
representatives acting on its certified test equilibria. At the other end
the orthogonal group of each of the nine six-body representatives acts on
the six choices of the marked body. For each orbit it records the
orthogonal stabilizer of the mark and the cycle type of its reflection.
The test sends a candidate primary to the gauge anchor, tests both
determinants and requires inclusion of the transformed coordinates in
the original uniqueness boxes. A verified symmetry sends a restricted
equilibrium to exactly one of the exhaustive disjoint root boxes.
Equivariant uniqueness and compactness preserve the stabilizers for
small positive $t$. For a marked shape with rotational stabilizer of
order $r$, its labelled oriented weight is $120/r$ with reflection and
$240/r$ without reflection.

\begin{theorem}
For some real algebraic $t_0\in(0,1)$ a noncollinear central
configuration with exactly bodies 5 and 6 on a reflection axis is
degenerate in the reflection-preserving shape equations. It is possible
to choose the configuration to have a unique mirror axis.
\end{theorem}
\begin{proof}
Let $g$ be a plane reflection followed by relabelling $(12)(34)$,
acting on the space of labelled oriented shapes. The reflection is
well-defined on the similarity quotient. With bodies 5 and 6 fixed
at $(1,0)$ and $(0,0)$ it is simply complex conjugation. A $g$-fixed
configuration must pair 1 with 2 and 3 with 4, and neither pair can
meet the axis without a collision. The fixed shape space is consequently
a smooth open subset of a four-dimensional linear subspace.
The invariant potential has vanishing normal derivative at a fixed
point. Its critical points in this subspace are exactly the $g$-fixed
physical central configurations. Restricted nondegeneracy gives local
continuation by the implicit function theorem.

There are 15 double transpositions of the five heavy labels. An ordinary
two-axial marked class of weight 120 contributes $120/15=8$ solutions
to this particular fixed space. The pentagon class contributes
$24\cdot5/15=8$. This follows by conjugacy in $S_5$; no other class
contributes. Thus the two counts are $8(13+1)=112$ and
$8(11+1)=96$. On every closed positive mass interval all the normalized
physical solutions form a compact family~\cite{Small}. Compactness is
unchanged on the fixed subspace and in our chart, since the two
distinguished bodies remain distinct. If the restricted Jacobian were
invertible at every fixed solution between the two endpoint regimes,
the number would be locally constant throughout, a contradiction.
The two endpoint counts are valid on open mass intervals by full
nondegeneracy. The regular pentagon is nonsingular on $(0,1)$ by
Meyer and Schmidt~\cite{MS} and by the accompanying exact Hessian
calculation. The remaining shape has just one axis. All equations,
distance and mass inequalities, and the vanishing restricted Hessian
determinant are semialgebraic over $\mathbb Q$. By real closed field
transfer there is a witness with algebraic mass and coordinates.
\end{proof}
The two four-axial pitchforks in the companion manuscript are geometrically
distinct from the existential degeneracy just proved, regardless of whether some of
their parameter values also support other singular shapes. Their local
changes remove a total of $2\cdot120\cdot2=480$ labelled oriented solutions
as $t$ increases through the two disjoint neighbourhoods. The agreement
with $4104-3624$ does not recover the intervening diagram of real solutions.
In particular, neither event removes a two-axial class, while the endpoint
counts for that class differ.

\section{Three saddle-nodes with two bodies on the reflection axis}
\label{sec:two-axial-folds}
Take one unit mass at $(0,0)$ and the body of mass $t$ at $(1,0)$.
The four other unit masses occupy $(a,\pm b)$ and $(c,\pm d)$, with $b,d>0$.
This gauge fixes translation, rotation and scale using two distinct labelled
bodies. We use Newtonian acceleration with gravitational constant one.
For $h(x,y)=(x^2+y^2)^{-3/2}$, put
\[
 (z_1,z_2,w_1,w_2,s,v)=
 (h(a,b),h(c,d),h(a-1,b),h(c-1,d),h(a-c,b-d),h(a-c,b+d)).
\]
Define $A=t+2az_1+2cz_2$ and $L=A+1-2(a-1)w_1-2(c-1)w_2$.
The symmetric equations $E=0$ have four components,
\begin{align*}
 E_1 &= a(L-z_1)-A+t(1-a)w_1+(c-a)(s+v),\\
 E_2 &= b(L-z_1-tw_1)+(d-b)s-(d+b)v-1/(4b^2),\\
 E_3 &= c(L-z_2)-A+t(1-c)w_2+(a-c)(s+v),\\
 E_4 &= d(L-z_2-tw_2)+(b-d)s-(b+d)v-1/(4d^2).
\end{align*}
Let $J=D_{(a,b,c,d)}E$ and let $K$ denote the complementary four-dimensional
derivative in the full planar chart, defined below.

\paragraph{Theorem.}
There are three distinct, real algebraic mass parameters with the following
properties. Each parameter differs from the corresponding table center by
less than $10^{-11}$. At the corresponding collision-free planar central
configuration, the normalized solution set is a smooth local curve with
coordinate $a$. The mass restricted to this curve satisfies $t'(a)=0$,
with second derivative enclosed below. Near either local maximum there are
two local oriented labelled configurations below the critical mass, one at
the critical mass, and none above it. At the local minimum the sides are
reversed. These counts apply in a neighborhood of the tabulated shape.
\begin{center}
\begin{tabular}{lll}
Mass center & Enclosure of $t''(a)$ & Enclosure of $\det K$\\
$0.00095963933728700494$ & $[-48.105221,-48.104660]$ & $[31.592665,31.592667]$\\
$0.1150143634630535$ & $[-6.566172,-6.566163]$ & $[192.066653,192.066659]$\\
$0.99770185629795693$ & $[26.721654,26.721813]$ & $[417511.421776,417511.432504]$
\end{tabular}
\end{center}

\paragraph{Full planar equations.}
Let $a_i$ denote the acceleration of body $i$, and write $q_0=(0,0)$,
$q_1=(1,0)$ for the pinned bodies. Use $R(x,y)=(-y,x)$,
$L=(a_0)_x-(a_1)_x$, and $W=(a_1)_y-(a_0)_y$. The residual
$a_i-a_0+Lq_i-WRq_i$ vanishes identically at $i=0,1$. If it vanishes
at the four free bodies as well, summing forces and taking total torque
about the center of mass shows that $W$ times the moment of inertia
vanishes. Thus $W=0$, and these are exactly the central-configuration
equations. At the symmetric shapes $L$ agrees with the formula above.
Reflection across the axis combined with transposition of the two labels
in each pair splits the derivative into two blocks. The even block is $J$.
In the odd block the displacements within each pair take the form
$(u,v)$ and $(-u,v)$. The four components of $K$ are the linearized
two upper-body residuals under these variations. To evaluate them the
verifier uses the Newtonian pair derivative
\[
 m_j\bigl(|D|^{-3}I-3|D|^{-5}DD^{\mathsf T}\bigr)
 (\delta q_j-\delta q_i),\qquad D=q_j-q_i.
\]
The nonzero bounds on $\det K$ ensure, by the implicit function theorem
and reflection symmetry, that all nearby full planar solutions are on
the symmetric locus.

\paragraph{Interval proof.}
The accompanying verifier \texttt{two\_axial\_three\_folds.py} uses only the
Python standard library. It evaluates the augmented equations
$G=(E,\det J)$ and their derivatives by second-order automatic
differentiation with dyadic interval endpoints and outward-rounded integer
square roots. Each five-coordinate box has radius $r=10^{-11}$ about the
rational center specified in the script. Its rational matrix $B$ is a
numerical preconditioner checked by the exact inequalities
\[
 \|I-BDG(X)\|_\infty\leq 10^{-6},\qquad
 \|BG(x_0)\|_\infty/r\leq 257\cdot10^{-6}.
\]
The contraction theorem supplies a unique zero in each box.
In each box the verifier also proves that
$H=D_{(b,c,d,t)}E$ is invertible. Along the curve parameterized by $a$,
write $u=(1,-H^{-1}E_a)$. The augmented condition $\det J=0$ implies $t'=0$
by Cramer's rule. The verifier encloses $t''$ as the last component of
$-H^{-1}D^2E[u,u]$. The displayed bounds give the three local fold
classifications. Finally the box-unique zeros are real algebraic:
all equations and their derivatives on the positive-distance chart are
semialgebraic over $\mathbb Q$, as are the rational isolating boxes.
The mass can therefore be picked out uniquely over the real algebraic
numbers.

\section{A two-axial pitchfork that creates two asymmetric configurations}
\label{sec:two-axial-pitch}
There is a second symmetry type in which the asymmetric branches appear on
the higher-mass side. For the mass line $(1,1,1,1,1,t)$, let
\[
 t_3^0=0.2972453330040611210936935514361686 .
\]
For some $|t_3-t_3^0|\le 10^{-15}$, there is a two-axial reflection-symmetric
configuration $Q_3$ that is locally unique at that mass. In a neighbourhood of
$Q_3$ in labelled, oriented shape space, the local count is one for
$t\le t_3$ and three for $t>t_3$, with $t$ sufficiently close to $t_3$.
The additional solutions are mirror related after relabelling and have no
physical reflection symmetry. Near this configuration the count stays at most
three when all six masses are perturbed independently.

Number particles from zero and divide the masses by five. Write
$F_i=q_i-\sum_{j\ne i}m_j(q_i-q_j)/|q_i-q_j|^3$, fix the angular coordinate
by $q_{3y}=0$, $q_{3x}>0$, and drop the residual $F_{3y}$. The lost
equation follows from the zero-torque identity. In this chart the involution
$S$ combines horizontal reflection with $(0\;2)(1\;4)$.
In the even subspace write
\[
 (q_0,q_1,q_2,q_3,q_4,q_5)
 =((a,b),(c,d),(a,-b),(e,0),(c,-d),(f,0)).
\]
The coordinates $(a,b,c,d,e,f)$ of $Q_3$ are approximately
\[
(-.544303,.390337,.213614,.711277,.371489,.975249).
\]
An odd variation has positions
\[
((v_0,v_1),(v_2,v_3),(-v_0,v_1),(0,0),(-v_2,v_3),(0,v_4)).
\]
The normalized null vector has $v_4=1$. The collision-free chart, the
distinctness of the anchor radius from the other heavy-particle radii, and
the separation of all labelled positions fix the only possible physical
reflection to be $S$ in this neighbourhood.

Here is the interval certificate in this chart.
Let $J_+,J_-$ denote the even and odd blocks of the force Jacobian.
The file \texttt{verify\_two\_axis.py} reads \texttt{two-axis.json}.
In a rational cube of radius $10^{-15}$ it certifies the unique solution
of $F_+=0$, $J_-v=0$, $v_4=1$. For a square rational preconditioner $R$,
the residual is bounded by $\|RG(x_0)\|_\infty<10^{-25}$ at the centre,
and $\|I-RDG(X)\|_\infty<4\cdot10^{-12}$ over the cube.
The even Jacobian and
$M=\left(\begin{smallmatrix}J_-&e_2\\ e_4^{\mathsf T}&0\end{smallmatrix}\right)$
are invertible throughout. The odd cokernel projection $\rho$ sends $y$ to
the last component of $M^{-1}(y,0)$.
Let $s$ denote the fifth odd coordinate of the position
displacement. On solving the other ten equations, write $g(s,t,\delta)$
for the remaining $\rho$-projected residual, where only the masses of
particles 1 and 4 are split by $1+\delta$ and $1-\delta$ before division
by five. The certified intervals are
\begin{align*}
 A=g_{st}&\in[-0.941439597,-0.941439596],\\
 C=g_{sss}&\in[2.336689820,2.336689821],\\
 B=g_\delta&\in[0.089844739,0.089844740].
\end{align*}
For completeness, if $w=-J_+^{-1}F_t$ and $u=-J_+^{-1}D^2F[v,v]$
are lifted to even positions, the three formulae evaluated are
$A=\rho[(\partial_t DF)v+D^2F[v,w]]$,
$C=\rho[D^3F[v,v,v]+3D^2F[v,u]]$ and $B=\rho F_\delta$.
They follow by implicit differentiation, since $\rho$ annihilates the
linear range corrections.

On the symmetric mass line $g$ is odd in $s$. Dividing by $s$ and applying
the implicit function theorem in $s^2$ gives the nontrivial branch
$t=t_3+\kappa s^2+O(s^4)$, where
$\kappa=-C/(6A)\in[.413673171,.413673172]$. Thus two symmetry-breaking
branches appear as $t$ increases. Weierstrass preparation of the reduced
force gives a monic cubic times a nonvanishing factor for all nearby
masses, which proves the asserted local bound. After removing the cubic's
quadratic term, the linear and constant coefficients are analytic local
coordinates in $(t,\delta)$ because $AB\ne0$ and $g_t=0$ at the tip.
Consequently these two physical parameters give a versal cusp.
Its fold arms satisfy
$|\delta|\sim K(t-t_3)^{3/2}$, with
$K=\sqrt{-8A^3/(9B^2C)}\in[6.270746461,6.270746462]$.

\section{What changes along the physical mass line}
The interval calculations in this manuscript and its companion identify
six distinct parameter values below one. The entries count labelled
oriented solutions in one neighbourhood of one marked configuration;
neighbourhoods related by the 120 permutations of the heavy labels are
disjoint in the labelled oriented space. At a fold a local maximum of the
mass loses two real solutions and a local minimum gains two. At a pitchfork
a mirror-symmetric solution persists while a pair of asymmetric solutions
is born or lost. All six fixed-mass critical points listed here are
isolated. The parameters in this summary are rounded, with rigorous
enclosures given in the corresponding sections and certificates.
\begin{center}
\begin{tabular}{rlcl}
\toprule
Mass ratio & Type of event & Bodies on axis & Local counts (below, at, above)\\
\midrule
$0.000959639337$ & fold & 2 & $2,1,0$\\
$0.089353633682$ & pitchfork & 4 & $3,1,1$\\
$0.115014363463$ & fold & 2 & $2,1,0$\\
$0.175165824530$ & pitchfork & 4 & $3,1,1$\\
$0.297245333004$ & pitchfork & 2 & $1,1,3$\\
$0.997701856298$ & fold & 2 & $0,1,2$\\
\bottomrule
\end{tabular}
\end{center}
The marked pentagon has rotational order five. At each tabulated
critical shape the rotational stabilizer is trivial: the points
are not the centre and vertices of a regular pentagon, as verified by
their coordinate enclosures. The only marked symmetry is the described
reflection. In the convention of the endpoint table, two downward and one
upward fold give a net loss of two ordinary two-axial classes, and two
downward and one upward pitchfork give a net loss of one chiral pair.
Consequently their total net loss of labelled oriented solutions is
$2(120)+1(240)=480$. This accounts arithmetically for the endpoint
difference. Further births and losses can have zero net contribution;
the calculations do not give global counts between the parameter boxes.

The two four-axial pitchforks, at masses near 0.08935 and 0.17517, are
proved in the companion manuscript \emph{Two symmetry-breaking bifurcations
of planar Newtonian six-body central configurations} (Chatgpt, 2026).
A separate certificate included with the present manuscript proves the
third pitchfork using the same physical chart and an independent check
in the two-axial gauge. Dias and Pan~\cite{DP} studied generic finiteness
in the cross-shaped symmetry class. The asymmetric solutions reached
through the four-axial pitchforks leave that class.

\section{Recovering the sixth mass from a marked physical shape}
The following calculation also applies when the five primary masses
$a_i>0$ are unequal and known. Put the sixth body at the origin, write
$x_i=q_i-q_6$, and set
\[
s_i=|x_i|^{-3},\quad
H_i=\sum_{j\ne i,\ j\le5}a_j(x_j-x_i)|x_j-x_i|^{-3},
\quad A=\sum_{i=1}^5 a_i s_i x_i .
\]
Here $A$ is the acceleration of the sixth body. Its subtraction from the
five primary accelerations gives
$H_i-A=(t s_i-\lambda)x_i$. Set $K_i=(H_i-A)\cdot x_i/|x_i|^2$.
The five vectors $H_i-A$ must be parallel to their $x_i$, and the
points $(s_i,K_i)$ must have a common affine relation $K_i=ts_i-\lambda$.
These two conditions are also sufficient when $t>0$, since the
pairwise central forces have zero mass-weighted sum. Their virial
identity gives $\lambda>0$. Define
\[
 D=5\sum s_i^2-(\sum s_i)^2=\sum_{i<j}(s_i-s_j)^2,\qquad
 N=5\sum s_i K_i-(\sum s_i)(\sum K_i).
\]
Whenever two radii differ, $D>0$ and $t=N/D$. A physical shape can
support at most one sixth mass there. This real sum-of-squares
deduction should not be used to remove a complex component merely because
$D$ vanishes on it.

For completeness suppose all radii have the same value $R$ and a
positive mass works. Write $M=\sum a_i$, $c=M^{-1}\sum a_ix_i$,
$s=R^{-3}$ and $\gamma=\lambda-ts$. Summing the primary equations
shows that $A=\gamma c=sMc$. Thus the primaries by themselves form
a cocircular five-body central configuration of multiplier $\gamma>0$.
If $c\ne0$, necessarily $R=(M/\gamma)^{1/3}$, excluded by
Yu and Zhu, Proposition 2~\cite{YZ}. Hence $c=0$. Wang's centred
cocircular classification~\cite{Wang} gives a regular pentagon with
equal primary masses. The converse follows directly by symmetry:
any positive mass at its centre works. This inverse uniqueness
identifies the physical locus on which the mass cannot be recovered.
It leaves open whether a fixed mass can have infinitely many distinct
shapes.

% Standalone section, or include in any LaTeX article using amsmath.
\section{The central pentagon branch}

Put five labelled unit masses at the vertices of a regular pentagon of
circumradius one and a mass $t>0$ at its centre. Meyer and Schmidt computed
the positive degeneracy of this branch in 1988, with approximate value
$6.478$ (\emph{SIAM J. Math. Anal.} 19, 1295--1313, Appendices A--B,
doi:10.1137/0519094). We give a Cartesian verification, an exact algebraic
specification of the parameter and the two indices on either side.

Write $a=2\sin(\pi/5)$, $b=2\sin(2\pi/5)=(3-a^2)a$ and
$h=a^{-1}+b^{-1}$. Then $a^4-5a^2+5=0$ and $1.17<a<1.18$.
The multiplier at radius one is $t+h$.
In Cartesian coordinates let $H$ be the full Hessian of the positive
Newtonian potential plus $(t+h)\sum_i m_i|q_i|^2/2$. Fix the first
vertex on the positive horizontal axis and delete its vertical coordinate.
The resulting principal matrix $H_*$ is the restriction to a complement
of the rotation tangent. Translation and dilation contribute only positive
directions for the mass metric. Thus the number of negative eigenvalues
of $H_*$ is the physical shape Morse index.

There is a nonzero constant $c$ independent of $t$ such that
\[
 \det H_* = c\,t^2(t+h)^3(t-t_-)^2(t-t_c)^2,
\]
where
\[
 t_-=-\frac{24a^3}{605}-\frac{37a}{242},\qquad
 t_c=\frac{-345208a^3-320500a^2+1180980a+1216925}{247215}.
\]
Here $t_-<0$. The unique positive degeneracy of this branch is $t_c$,
the zero of the following irreducible polynomial in the interval from 6.478190058 to 6.478190060,
\[
 6674805t^4-44892900t^3+8565750t^2+13531500t+2057501 .
\]
At this parameter $H_*$ has rank nine and nullity two. The Morse index is
zero for $0<t<t_c$ and two for $t>t_c$. Both zero eigenvalues cross simply
with negative derivative when $t$ increases. In particular the count
transition known to occur at some $0<t<1$ on the five-equal-mass ray occurs
on another branch.

For a reproducible proof, construct each pair block
$m_im_j(3dd^{\mathsf T}/|d|^5-I/|d|^3)$ for $d=q_i-q_j$, add it to
the two diagonal blocks of $H$ and subtract it from the two off-diagonal
blocks. The pentagon coordinates, all three distances $a,b,1$ and all
coefficients belong to $\mathbb Q(a)$. In the sliced matrix the nine outer
rows have degree at most one in $t$, and the two centre rows have degree
at most two. Fourteen distinct integer parameter evaluations therefore
recover its determinant identically. The accompanying script performs
this interpolation and checks the displayed factorization exactly.
Polynomial elimination, irreducibility and a Sturm count give the stated
isolating interval. Exact interval evaluation in the specified embedding
determines the signs of symmetric elimination pivots at $t=1$ and $7$.
The indices are constant on each nonsingular interval. At $t_c$ the
order two of the determinant equals the nullity, so both vanishing
eigenvalues of this real analytic symmetric matrix have order one,
and their signs of crossing follow from the index change. The projection
of the nullspace onto the central particle's two displacement components
is an isomorphism: adjoining those projection rows gives a rank-eleven
matrix. This identifies the kernel as the planar displacement mode for
the combined rotation and cyclic relabelling of the pentagon.


\subsection{The critical fibre and its local multiplicity}
Use the invariant affine slice perpendicular to the geometry's rotation
tangent $\tau=(Jq_0,\ldots,Jq_5)$ in the ordinary Cartesian inner
product. Write $u$ for nine free outer-ring displacements and $(x,y)$
for central displacement. The $9\times9$ ring Hessian $A$ at $t_c$ is
invertible by the accompanying exact calculation. The implicit function
theorem solves the nine ring stationarity equations for $u=u(x,y)$,
and we write $\Psi$ for the energy restricted to this graph. The slice,
the equations defining the graph, and its unique solution respect the
dihedral symmetry $D_5$ (rotation followed by cyclic relabelling, and
reflection followed by reversal of vertex order). A slice critical
point solves the full system: the infinitesimal rotation remains
transverse to the slice nearby and the energy is rotation-invariant.

Write $v$ for the tangent to the stationary graph as $(x,y)=(s,0)$ varies,
so its ring components are $-A^{-1}\Phi_{ux}$ and its central components
are $(1,0)$. With $b=\Phi_{u v v}$ at the reference point, the quadratic
coefficient of $u(s,0)$ is $-A^{-1}b/2$. Substitution of this quadratic
Taylor approximation into the potential gives its expansion through fifth
order: its error in $u$ has size $O(s^3)$ and the first variation in $u$
on the stationary graph is zero. The script expands each inverse
distance by the binomial series, and obtains coefficients in $\mathbb Q(a)$.
In the specified real embedding it verifies the rational bounds
\[
 -69.721735<c_4<-69.721734,\qquad 2167.520399<c_5<2167.520400,
\]
where $\Psi(s,0)=\Psi(0,0)+c_4s^4+c_5s^5+O(s^6)$.

Complexify by putting $z=x+iy$ and $w=x-iy$. The branches of all inverse
distances are chosen to agree with positive physical distances at the
base point. Rotation gives weights $1,-1$ modulo five, and reflection
interchanges the two coordinates. The vanishing of the linear and
quadratic terms, and the dihedral invariance of the remaining terms,
therefore give
\[
 \Psi(z,w)=\Psi(0,0)+c_4z^2w^2+\tfrac12 c_5(z^5+w^5)
                 +O((|z|+|w|)^6).
\]
Let $P=\Psi_z$, $Q=\Psi_w$. The function $zP-wQ$ changes sign under
each of the five reflections and vanishes on each complex reflection
line $w=\zeta z$, $\zeta^5=1$. Dividing its analytic germ by the five
distinct linear factors gives
\[
 zP-wQ=(z^5-w^5)E(z,w),\qquad E(0,0)=5c_5/2\ne0.
\]
On each line $P(z,\zeta z)=2c_4\zeta^2z^3+O(z^4)$. Hence $P=Q=0$
has only the origin nearby, even over the complex numbers. To compute its
local intersection multiplicity $i_0(P,Q)$ use the product rule for local
plane-curve intersection numbers. Since $P(z,0)=(5c_5/2)z^4+O(z^5)$,
\[
 15=\sum_{\zeta^5=1}i_0(P,w-\zeta z)
    =i_0(P,zP-wQ)=i_0(P,wQ)=4+i_0(P,Q).
\]
The multiplicity at the critical mass is exactly eleven. Eliminating nine
variables by the implicit function theorem preserves this local multiplicity.
Conservation of number for isolated holomorphic zeros gives a fixed
neighbourhood of the reference shape in which sufficiently small
independent variations of the six masses have total complex multiplicity
eleven; there are at most eleven distinct real labelled shapes therein.
One may fix the positive multiplier at $t_c+h$ when varying all masses:
homogeneity uniquely fixes the scale for every nearby physical shape.
For $t\ne t_c$ the invertible sliced Hessian similarly gives local
multiplicity one throughout the central pentagon branch.

For comparison with the classical real bifurcation, in the symmetric family
with $t$ near $t_c$, the reduced energy has quadratic part $p(t)(x^2+y^2)$,
with $p(t_c)=0$ and $p'(t_c)<0$ by the index crossing and the Schur
complement. The same reflection-line factorization holds for nearby
$t$, since $E(0,0)$ remains nonzero. Along any of the ten real reflection
rays the radial critical equation, divided by its positive radius $r$,
is $2p(t)+4c_4(t)r^2+O(r^3)=0$. The left-hand side has negative
radial derivative for all sufficiently small positive $r$, uniformly
in $t$, while $c_4(t_c)<0$. Thus there are ten additional noncentral
solutions just below $t_c$ and none just above it, in the labelled
rotation slice. This agrees with the dihedral bifurcation described by
Meyer and Schmidt (1988, Sec.\ 5 and Appendix C). All the multiplicities
and counts here concern one neighbourhood in shape space.

\section{Verification and provenance}
The accompanying archive supplies the exact specialization tests, interval
verifiers, input data, Lean source and transcripts. The diagram sign
tables, cubic nonvanishing tests and root-disk inequalities use Lean kernel evaluation on the
dyadic centres. Lean verifies the shape polynomial products, while
Rouch\'e's theorem, the resultant argument, and the identification of
these algebraic tests with the singular diagram are supplied by the
written proof. For pitchforks the formal proof treats the
root count of the reduced scalar equation with stated analytic hypotheses.
The bridge from the Newtonian equations to its derivatives and cubic
coefficients uses the interval code and the paper. For folds the interval
verification is in Python. The pentagon computations use exact arithmetic
in a quartic number field; Lean independently checks the specified
polynomial and coefficient identities. The endpoint counts depend on the
published completeness theorem and the previous small-mass research
manuscript. These statements have different verification scopes.
All original manuscripts and accompanying documentation are licensed
under CC BY 4.0, and original verification code under MIT.
\begin{thebibliography}{9}
\bibitem{Tomar} T. S. Tomar, \emph{Smale's sixth problem for six bodies in the plane:
finiteness of central configurations for generic masses, with a proof formalized in Lean},
research preprint (2026), \url{https://hexagonmath.org/2611.00377v1}.
\bibitem{MZ} M. Moczurad and P. Zgliczy\'nski, \emph{Central configurations in planar
$n$-body problem with equal masses for $n=5,6,7$}, Celest. Mech. Dyn. Astron.
131 (2019), article 46. \url{https://doi.org/10.1007/s10569-019-9920-6}.
\bibitem{Small} Chatgpt, \emph{Exactly 4104 labelled oriented planar central
configurations for five nearly equal masses and one light mass}, research manuscript,
with input enclosures and code (2026).
\href{https://github.com/rongyulin/physics-open-problems-AI/tree/523559700ce3c996afe6b08f4744309cfb631d3f/problems/planar-central-configurations}{Archived release and companion manuscripts}.
\bibitem{MS} K. R. Meyer and D. S. Schmidt, \emph{Bifurcations of relative equilibria
in the $N$-body and Kirchhoff problems}, SIAM J. Math. Anal. 19 (1988), 1295--1313.
\url{https://doi.org/10.1137/0519094}.
\bibitem{DP} T. Dias and B.-Y. Pan, \emph{Generic finiteness for a class of
symmetric planar central configurations of the six-body problem and the six-vortex
problem}, J. Dynam. Differential Equations 32 (2020), 1579--1602.
\url{https://doi.org/10.1007/s10884-019-09765-y}.
\bibitem{YZ} X. Yu and S. Zhu, \emph{Classification of stacked central configurations in
$\mathbb R^3$}, \url{https://arxiv.org/abs/1812.00200}.
\bibitem{Wang} Z. Wang, \emph{On the centered co-circular central configurations for
the $n$-body problem}, J. Differential Equations 361 (2023), 376--390.
\url{https://doi.org/10.1016/j.jde.2023.02.066}.
\end{thebibliography}
\end{document}
